% FILE: 00_project_identity.tex
% Project: adversarial
% Thesis Role: adversarial-testing-and-hostile-assumption-verification

\section{Project Identity: Adversarial Testing and Hostile Assumption Verification}
\label{sec:adversarial-identity}

\subsection{Purpose}
\label{sec:adversarial-identity-purpose}

The \texttt{adversarial/} directory is the red-team and hostile-assumption-verification layer
of the process-intelligence research foundry. Its governing mandate is stated in
\texttt{RED\_TEAM\_SCOPE.md}: ``Every claim in this research corpus must survive adversarial
scrutiny before it becomes a board artifact or a manufacturing receipt.'' This mandate
operationalizes the Van der Aalst Constitution requirement for negative testing---the
principle that model-versus-log mismatch is not a discrepancy but a first-class defect.

The adversarial corpus does not simply catalog attack scenarios in the abstract. It provides
executable tests, cryptographic challenge mechanisms, and a structured finding taxonomy that
feeds directly into gap remediation sequencing and downstream workflow authorization. A claim
that cannot be challenged is not a research claim; it is an assertion. The adversarial layer
exists to enforce the difference.

The corpus operates under the Chicago TDD hostile assumptions doctrine, which requires that
declared manufacturing pipelines be distinguished from actual runtime processes, that receipts
be verified against cryptographic evidence rather than accepted on declaration, and that proof
gates be shown to reject non-conforming execution---not merely to pass conforming execution.

\subsection{Repository Surface}
\label{sec:adversarial-identity-surface}

The adversarial corpus spans three interlocking mechanisms. The first is a scoped review
charter (\texttt{RED\_TEAM\_SCOPE.md}) defining seven challenge categories: unsupported claims
in \texttt{doctrine/}; inflated paper classifications (\texttt{COVERED\_BY\_TYPE} without
actual types); board claims without evidence paths; lifecycle phases that stop at observation
without repair actuation; comparison claims without matrix evidence; crosswalk claims without
loss policy specification; and M\&A claims without a diligence path from target OCEL log to
conformance report to board projection.

The second mechanism is an executable simulation engine
(\texttt{simulate\_laundering\_spoofing.py}) implementing ten attack scenarios against the
\texttt{IngestionBoundary} class. The third is a standards audit script
(\texttt{validate\_standards.py}) that validates placement files for XES, OCEL 2.0, BPMN,
POWL, Declare, and OCPQ.

Supporting the three core files are adversarial documents distributed across the wider
foundry: \texttt{standards/ocel-adversarial-gravity.md} (ontological vulnerability analysis),
\texttt{sources/papers/adversarial-canon.md} (Sybil-attack failures in classical discovery
algorithms), and \texttt{sources/wasm4pm-compat/adversarial-type-law.md} (witness lattice
forgeability and the Chatman Constraint).

\subsection{Primary Research Role}
\label{sec:adversarial-identity-role}

Within the dissertation, the adversarial corpus serves as the falsification layer for all
positive claims. Every ALIVE gate criterion in the research program is defined against a
background of adversarial pressure: the conformance claim must survive signature forgery, the
process model claim must survive event deletion and replay, the board claim must survive the
challenge ``show me the diligence path.''

The adversarial corpus also makes a positive research contribution: it demonstrates that
standard public process mining formats (XES, OCEL 2.0) are structurally vulnerable to silent
tampering, timeline manipulation, and phantom-object injection, and that the
\texttt{wasm4pm} Reality Receipt mechanism addresses these vulnerabilities at the ingestion
boundary through cryptographic hash chaining, HMAC-SHA256 signature verification, temporal
monotonicity gates, and impossible-velocity detection.

\subsection{Key Artifacts}
\label{sec:adversarial-identity-artifacts}

The key manufactured artifacts are:
\begin{itemize}
  \item \texttt{adversarial/RED\_TEAM\_SCOPE.md} --- the charter for all adversarial activity,
        defining seven challenge categories and the output format for findings;
  \item \texttt{adversarial/RED\_TEAM\_FINDINGS\_001.md} and
        \texttt{adversarial/RED\_TEAM\_FINDINGS\_002.md} --- ten structured findings (3
        CRITICAL, 3 MAJOR open, 2 self-correcting, 3 MINOR) with severity, category, and
        remediation paths;
  \item \texttt{adversarial/simulate\_laundering\_spoofing.py} --- ten executable scenarios
        with Python \texttt{assert} statements verifying ACCEPTED and REJECTED verdicts;
  \item \texttt{adversarial/validate\_standards.py} --- standards audit script covering six
        formats with placeholder, broken-link, JSON validity, and LossReport checks;
  \item \texttt{checkpoints/PROCESS\_INTELLIGENCE\_ADVERSARIAL\_V30.1.1\_OMEGA.md} --- the
        ALIVE OMEGA checkpoint issued after a twenty-agent adversarial swarm;
  \item \texttt{audits/adversarial\_audit\_v30.1.1.md} --- VERIFIED/PASS verdict with Blake3
        checksum and Ed25519 signature over 567 commits across 15 directories.
\end{itemize}
